\clearpage
\VerdierPhysicalPageAnchor{0169}%
\begingroup\small\setlength{\parskip}{0pt}\setlength{\emergencystretch}{2.5em}
\VerdierSemanticAnchor{ch4-4-2-6-third-exact-sequence}
\[
0\longrightarrow
E\!\begin{bmatrix}n\\ f_2f_1,f_3,f_4\end{bmatrix}
\longrightarrow
E\!\begin{bmatrix}n\\ f_2,f_3,f_4\end{bmatrix}
\longrightarrow
Z\!\begin{bmatrix}n\\ f_2,f_3\end{bmatrix}
 \Big/
Z\!\begin{bmatrix}n\\ f_2f_1,f_3\end{bmatrix}
\longrightarrow0.
\]
Soit alors $f_i$, $i\in\mathbb Z$, une suite illimitée de morphismes
composables de $\mathcal J$. On déduit de ce qui précède et de la définition
des différentielles, une suite exacte illimitée :
\VerdierSemanticAnchor{ch4-4-2-6-unbounded-exact-sequence}
\[
\begin{aligned}
\cdots\longrightarrow
E\!\begin{bmatrix}n-1\\ f_{i+3},f_{i+4},f_{i+5}\end{bmatrix}
&\xrightarrow{d^{n-1}}
E\!\begin{bmatrix}n\\ f_{i+1},f_{i+2},f_{i+3}\end{bmatrix}
\longrightarrow
E\!\begin{bmatrix}n\\ f_{i+1},f_{i+3}f_{i+2},f_{i+4}\end{bmatrix}
\longrightarrow{}\\[-1pt]
&\longrightarrow
E\!\begin{bmatrix}n\\ f_{i+2},f_{i+3},f_{i+4}\end{bmatrix}
\xrightarrow{d^n}
E\!\begin{bmatrix}n+1\\ f_i,f_{i+1},f_{i+2}\end{bmatrix}
\longrightarrow\cdots.
\end{aligned}
\]
Le lecteur curieux élucidera les rapports de cette suite exacte avec la
suite exacte à 5 termes de bas degré des suites spectrales classiques.

\VerdierSemanticAnchor{ch4-4-2-7}
\medskip\noindent\textbf{4.2.7.}
Tout diagramme commutatif de $\mathcal A$ :
\[
\begin{tikzcd}[column sep=4em,row sep=2.7em]
&Y\arrow[dl]\arrow[d,"\varphi"]\arrow[dr,"\psi"]&\\
X'\arrow[r,"\varphi'"']&X\arrow[r,"\eta"']&X'',
\end{tikzcd}
\]
dans lequel la ligne est exacte, définit par l'intermédiaire de $\eta$ un
isomorphisme ([1] p.~316, Lemme 1.1) :
\VerdierSemanticAnchor{ch4-4-2-7-image-isomorphism}
\[
\operatorname{Im}(\varphi)\big/\operatorname{Im}(\varphi')
\simeq\operatorname{Im}(\psi).
\]
Soit
$\omega_3=(\mathord\bullet\xrightarrow{f_1}\mathord\bullet
\xrightarrow{f_2}\mathord\bullet\xrightarrow{f_3}\mathord\bullet)$
un objet de $\mathcal{F}\ell_3(\mathcal J)$. En appliquant la remarque
précédente au diagramme :
\VerdierSemanticAnchor{ch4-4-2-7-first-diagram}
\[
\begin{tikzcd}[column sep=4.6em,row sep=2.8em]
&H^n(f_2f_1)\arrow[d]\arrow[dr]&\\
H^{n-1}(f_3)\arrow[ur,"\delta"]\arrow[r,"\delta"']&
H^n(f_2)\arrow[r]&H^n(f_3f_2),
\end{tikzcd}
\]
on obtient un isomorphisme :
\VerdierSemanticAnchor{ch4-4-2-7-1}
\[
E\!\begin{bmatrix}n\\ f_1,f_2,f_3\end{bmatrix}
\xrightarrow{\ \sim\ }
\operatorname{Im}\bigl(H^n(f_2f_1)\longrightarrow H^n(f_3f_2)\bigr).
\tag{4.2.7.1}
\]
\endgroup
